
\section{Reproduction}
\label{app:reproduction}

This appendix describes how to reproduce the results of the thesis.
The source code, the proofs, and the scripts are in the repository of the thesis; the commit used for the measurements is \EvalCommit.

\subsection{Requirements}

The verifier needs Python~3.11 or newer (the measurements used Python~\EvalPython) and the library \texttt{lark}.
The external prover is Vampire (the measurements used \EvalVampire).
Vampire is called through the environment variables \texttt{EXTERNAL\_ATP\_BINARY} (path of the binary) or \texttt{EXTERNAL\_ATP\_COMMAND} (complete command), and is found in the search path by default.
Under Windows, the verifier is run in WSL.

\begin{verbatim}
python3 -m venv .venv && source .venv/bin/activate
python -m pip install -r requirements.txt -r requirements-dev.txt
export EXTERNAL_ATP_BINARY=/usr/local/bin/vampire
\end{verbatim}

\subsection{Checking Proofs}

A single proof is checked by passing the proof file; the problem file is looked up in the subdirectory \texttt{Problems} next to the proof, or given with \texttt{-p}.
Without a proof argument, the verifier checks all proofs of a directory (default \texttt{ProoVer2026/PRV}; further sets are \texttt{ProoVer2026/Test}, \texttt{ProoVer2026/Extra}, and \texttt{ProoVer2026/Deprecated}) and prints a summary with the score.

\begin{verbatim}
python main.py ProoVer2026/PRV/PRV010+1.s
python main.py proof.s -p problem.p -t 30
python main.py -d ProoVer2026/Extra
python main.py --binary-verdict ProoVer2026/PRV/PRV010+1.s
\end{verbatim}

\begin{table}[tbp]
\centering
\small
\caption{Command-line options of the verifier.}
\label{tab:app-cli}
\begin{tabular}{@{}l P{0.65\linewidth}@{}}
\toprule
Option & Meaning \\
\midrule
\texttt{-p}, \texttt{--problem} & Problem file; by default \texttt{Problems/<name>.p} next to the proof. \\
\texttt{-t}, \texttt{--timeout} & Wall-clock limit per proof in seconds (default 30). \\
\texttt{-d}, \texttt{--dir} & Batch mode: directory with proofs; the expected results are chosen from the directory name (\texttt{-e} overrides). \\
\texttt{--prefix} & Batch mode: only check proofs whose file name starts with the prefix. \\
\texttt{--binary-verdict} & Print \texttt{VerifiedBad} instead of \texttt{Unknown} when a step cannot be decided. \\
\bottomrule
\end{tabular}
\end{table}

The last line of the output is the SZS status.
Table~\ref{tab:app-status} lists the statuses.

\begin{table}[tbp]
\centering
\small
\caption{Statuses printed by the verifier.}
\label{tab:app-status}
\begin{tabular}{@{}l P{0.65\linewidth}@{}}
\toprule
Status & Meaning \\
\midrule
\texttt{VerifiedGood} & Every check succeeded. \\
\texttt{VerifiedBad} & A definite defect was found, or the proof could not be processed. \\
\texttt{Unknown} & No defect was found, but a step could not be decided (prover without result, or no checker for the rule). \\
\texttt{Timeout} & The time limit was exceeded. \\
\bottomrule
\end{tabular}
\end{table}

\subsection{Tests and Evaluation}

The test suite is run with \texttt{pytest}; the coverage figures of Section~\ref{subsec:eval-modularity} come from the same run.
The evaluation driver runs the verifier in a fresh process for each proof and writes raw data to \texttt{results/eval} and the tables and macros of Chapter~\ref{sec:evaluation} to \texttt{thesis/generated}.

\begin{verbatim}
pytest
python scripts/eval_all.py --parts all --reps 3
python scripts/gen_appendix_tables.py
\end{verbatim}

\noindent The parts of the driver can be run separately with \texttt{--parts}; Table~\ref{tab:app-parts} lists them.
The run time depends on the machine; the full run needs the prover for all parts except \texttt{noatp}.

\begin{table}[tbp]
\centering
\small
\caption{Parts of the evaluation driver \texttt{scripts/eval\_all.py}.}
\label{tab:app-parts}
\begin{tabular}{@{}l P{0.70\linewidth}@{}}
\toprule
Part & Content \\
\midrule
\texttt{runs} & Verdicts, steps, prover calls, time, and memory per proof (Sections~\ref{subsec:eval-correctness} and~\ref{subsec:eval-runtime}). \\
\texttt{noatp} & The same without prover (ablation). \\
\texttt{internal} & Internal checkers against Vampire on single steps. \\
\texttt{mutation} & Mutations of accepted proofs with entailment check of the survivors. \\
\texttt{inputs} & Malformed inputs (robustness). \\
\texttt{extra} & Generated rule test suite. \\
\texttt{coverage} & Test coverage and lines of code per module. \\
\texttt{tables} & Regenerates tables, macros, and the figures from the raw data. \\
\bottomrule
\end{tabular}
\end{table}

\section{Example Reports}
\label{app:reports}

The report lists one block per check, and ends with the SZS status.
The following output is the (shortened) report for the invalid proof \texttt{PRV002+1}, in which the negated conjecture is not the negation of the conjecture; the blocks without finding are omitted and the formulas are abbreviated.

\begin{verbatim}
Checking proof: ProoVer2026/PRV/PRV002+1.s
Against problem: ProoVer2026/PRV/Problems/PRV002+1.p

[FAIL] Negated Conjecture Issues:
  - neg: formula is not the correct negation of parent
    conjecture; expected one of: [...], got: [...]

% SZS status VerifiedBad
\end{verbatim}

\section{Results per Proof}
\label{app:per-proof}

Tables~\ref{tab:app-prv}, \ref{tab:app-tst}, and~\ref{tab:app-ext} list the status and the run time of each proof of the three benchmark groups.
The time is the median over the repetitions of the whole process, including the start-up of the interpreter.
The column ``Result'' compares the status with the expected result; ``abstain'' denotes \texttt{Unknown} or \texttt{Timeout}, which score zero.

{\small
\begin{longtable}{@{}l l l l r r@{}}
\caption{Results for each proof of the ProoVer set (official proofs and samples).}\label{tab:app-prv}\\
\toprule
Proof & Expected & Status & Result & Time (s) & ATP calls \\
\midrule
\endfirsthead
\toprule
Proof & Expected & Status & Result & Time (s) & ATP calls \\
\midrule
\endhead
\bottomrule
\endlastfoot
\texttt{PRV000{+}1} & correct & VerifiedGood & ok & 1.69 & 1 \\
\texttt{PRV001{+}1} & correct & VerifiedGood & ok & 1.39 & 0 \\
\texttt{PRV002{+}1} & incorrect & VerifiedBad & ok & 1.58 & 1 \\
\texttt{PRV003{+}1} & incorrect & VerifiedBad & ok & 1.47 & 1 \\
\texttt{PRV004{+}1} & correct & VerifiedGood & ok & 1.40 & 1 \\
\texttt{PRV005{+}1} & correct & VerifiedGood & ok & 1.85 & 1 \\
\texttt{PRV006{+}1} & incorrect & VerifiedBad & ok & 1.66 & 1 \\
\texttt{PRV007{+}1} & incorrect & VerifiedBad & ok & 1.30 & 0 \\
\texttt{PRV008{+}1} & incorrect & VerifiedBad & ok & 1.27 & 0 \\
\texttt{PRV009{+}1} & incorrect & VerifiedBad & ok & 1.45 & 1 \\
\texttt{PRV010{+}1} & correct & VerifiedGood & ok & 1.42 & 0 \\
\texttt{PRV011{+}1} & correct & VerifiedGood & ok & 1.66 & 0 \\
\texttt{PRV012{+}1} & correct & VerifiedGood & ok & 1.44 & 0 \\
\texttt{PRV013{+}1} & correct & VerifiedGood & ok & 1.42 & 0 \\
\texttt{PRV014{+}1} & correct & VerifiedGood & ok & 1.73 & 1 \\
\texttt{PRV015{+}1} & correct & VerifiedGood & ok & 1.34 & 0 \\
\texttt{PRV016{+}1} & correct & VerifiedGood & ok & 1.10 & 0 \\
\texttt{PRV017{+}1} & correct & VerifiedGood & ok & 1.55 & 2 \\
\texttt{PRV018{+}1} & correct & VerifiedGood & ok & 1.42 & 0 \\
\texttt{PRV019{+}1} & correct & VerifiedGood & ok & 1.36 & 1 \\
\texttt{PRV020{+}1} & correct & VerifiedGood & ok & 1.23 & 0 \\
\texttt{PRV021{+}1} & correct & VerifiedGood & ok & 1.20 & 0 \\
\texttt{PRV022{+}1} & correct & VerifiedGood & ok & 1.29 & 2 \\
\texttt{PRV023{+}1} & incorrect & VerifiedBad & ok & 1.26 & 0 \\
\texttt{PRV024{+}1} & incorrect & VerifiedBad & ok & 1.28 & 0 \\
\texttt{PRV025{+}1} & incorrect & VerifiedBad & ok & 1.74 & 0 \\
\texttt{PRV026{+}1} & incorrect & VerifiedBad & ok & 2.12 & 1 \\
\texttt{PRV027{+}1} & correct & VerifiedGood & ok & 1.59 & 0 \\
\texttt{PRV028{+}1} & incorrect & VerifiedBad & ok & 1.53 & 1 \\
\texttt{PRV029{+}1} & correct & VerifiedGood & ok & 1.30 & 0 \\
\texttt{PRV030{+}1} & correct & VerifiedGood & ok & 1.34 & 1 \\
\texttt{PRV031{+}1} & incorrect & VerifiedBad & ok & 1.47 & 1 \\
\texttt{PRV032{+}1} & incorrect & VerifiedBad & ok & 1.24 & 0 \\
\texttt{PRV033{+}1} & correct & VerifiedGood & ok & 1.25 & 0 \\
\texttt{PRV034{+}1} & correct & VerifiedGood & ok & 1.43 & 0 \\
\texttt{PRV035{+}1} & incorrect & VerifiedBad & ok & 1.59 & 0 \\
\texttt{PRV036{+}1} & incorrect & VerifiedBad & ok & 1.42 & 0 \\
\texttt{PRV037{+}1} & incorrect & VerifiedBad & ok & 1.39 & 0 \\
\texttt{PRV038{+}1} & incorrect & VerifiedBad & ok & 1.34 & 1 \\
\texttt{PRV039{+}1} & incorrect & VerifiedBad & ok & 1.36 & 0 \\
\texttt{PRV040{+}1} & correct & VerifiedGood & ok & 1.41 & 1 \\
\texttt{PRV041{+}1} & correct & VerifiedGood & ok & 1.33 & 0 \\
\texttt{PRV042{+}1} & correct & VerifiedGood & ok & 1.69 & 1 \\
\texttt{PRV043{+}1} & correct & VerifiedGood & ok & 2.06 & 1 \\
\texttt{PRV044{+}1} & correct & VerifiedGood & ok & 1.90 & 1 \\
\texttt{PRV045{+}1} & correct & VerifiedGood & ok & 1.26 & 0 \\
\texttt{PRV046{+}1} & incorrect & VerifiedBad & ok & 1.42 & 0 \\
\texttt{PRV047{+}1} & incorrect & VerifiedBad & ok & 1.12 & 0 \\
\texttt{PRV048{+}1} & correct & VerifiedGood & ok & 1.34 & 0 \\
\texttt{PRV049{+}1} & incorrect & VerifiedBad & ok & 1.63 & 0 \\
\texttt{PRV050{+}1} & incorrect & VerifiedBad & ok & 1.84 & 0 \\
\texttt{PRV051{+}1} & incorrect & VerifiedBad & ok & 1.75 & 0 \\
\texttt{PRV052{+}1} & incorrect & VerifiedBad & ok & 1.59 & 0 \\
\texttt{PRV053{+}1} & incorrect & VerifiedBad & ok & 1.36 & 0 \\
\texttt{PRV054{+}1} & incorrect & VerifiedBad & ok & 1.50 & 0 \\
\texttt{PRV055{+}1} & correct & VerifiedGood & ok & 1.59 & 0 \\
\texttt{PRV056{+}1} & incorrect & VerifiedBad & ok & 1.70 & 1 \\
\texttt{PRV057{+}1} & incorrect & Error & ok & 1.89 & 0 \\
\texttt{PRV058{+}1} & incorrect & VerifiedBad & ok & 1.60 & 1 \\
\texttt{PRV059{+}1} & correct & VerifiedGood & ok & 1.53 & 1 \\
\texttt{PRV060{+}1} & correct & VerifiedGood & ok & 1.58 & 1 \\
\texttt{PRV061{+}1} & correct & VerifiedGood & ok & 1.86 & 3 \\
\texttt{PRV062{+}1} & correct & VerifiedGood & ok & 1.88 & 2 \\
\texttt{PRV063{+}1} & correct & VerifiedGood & ok & 2.00 & 1 \\
\texttt{PRV064{+}1} & correct & VerifiedGood & ok & 1.80 & 4 \\
\texttt{PRV065{+}1} & incorrect & VerifiedBad & ok & 1.77 & 1 \\
\texttt{PRV066{+}1} & incorrect & VerifiedBad & ok & 1.72 & 1 \\
\texttt{PRV067{+}1} & incorrect & VerifiedBad & ok & 1.38 & 0 \\
\texttt{PRV068{+}1} & incorrect & VerifiedBad & ok & 1.48 & 0 \\
\texttt{PRV069{+}1} & incorrect & VerifiedBad & ok & 1.52 & 1 \\
\texttt{PRV070{+}1} & correct & VerifiedGood & ok & 1.40 & 0 \\
\texttt{PRV071{+}1} & incorrect & VerifiedBad & ok & 4.24 & 1 \\
\texttt{PRV072{+}1} & incorrect & VerifiedBad & ok & 4.33 & 0 \\
\texttt{PRV073{+}1} & correct & VerifiedGood & ok & 4.27 & 1 \\
\texttt{PRV074{+}1} & incorrect & VerifiedBad & ok & 1.49 & 1 \\
\texttt{PRV075{+}1} & incorrect & VerifiedBad & ok & 1.51 & 1 \\
\texttt{PRV076{+}1} & correct & VerifiedGood & ok & 1.54 & 2 \\
\texttt{PRV077{+}1} & incorrect & VerifiedBad & ok & 1.41 & 1 \\
\texttt{PRV078{+}1} & incorrect & VerifiedBad & ok & 1.50 & 1 \\
\texttt{PRV079{+}1} & correct & VerifiedGood & ok & 1.62 & 1 \\
\texttt{PRV080{+}1} & incorrect & VerifiedBad & ok & 1.69 & 1 \\
\texttt{PRV081{+}1} & correct & VerifiedGood & ok & 1.61 & 1 \\
\texttt{PRV082{+}1} & correct & VerifiedGood & ok & 1.46 & 1 \\
\texttt{PRV083{+}1} & incorrect & VerifiedBad & ok & 1.41 & 0 \\
\texttt{PRV084{+}1} & incorrect & VerifiedBad & ok & 1.53 & 1 \\
\texttt{PRV085{+}1} & correct & VerifiedGood & ok & 1.54 & 1 \\
\texttt{PRV086{+}1} & incorrect & VerifiedBad & ok & 1.45 & 1 \\
\texttt{PRV087{+}1} & incorrect & VerifiedBad & ok & 1.72 & 1 \\
\texttt{PRV088{+}1} & correct & VerifiedGood & ok & 1.91 & 1 \\
\texttt{PRV089{+}1} & correct & VerifiedGood & ok & 1.59 & 0 \\
\texttt{PRV090{+}1} & incorrect & VerifiedBad & ok & 1.64 & 0 \\
\texttt{PRV091{+}1} & incorrect & Error & ok & 1.74 & 0 \\
\texttt{PRV092{+}1} & incorrect & VerifiedBad & ok & 2.32 & 1 \\
\texttt{PRV093{+}1} & incorrect & VerifiedBad & ok & 2.44 & 1 \\
\texttt{PRV094{+}1} & incorrect & VerifiedBad & ok & 2.31 & 1 \\
\texttt{PRV095{+}1} & incorrect & VerifiedBad & ok & 1.84 & 0 \\
\texttt{PRV096{+}1} & correct & VerifiedGood & ok & 1.79 & 0 \\
\texttt{PRV097{+}1} & incorrect & Unknown & abstain & 1.90 & 0 \\
\texttt{PRV098{+}1} & correct & VerifiedGood & ok & 2.10 & 0 \\
\texttt{PRV099{+}1} & incorrect & VerifiedBad & ok & 11.00 & 2 \\
\end{longtable}

}

{\small
\begin{longtable}{@{}l l l l r r@{}}
\caption{Results for each proof of the TST set (own test proofs).}\label{tab:app-tst}\\
\toprule
Proof & Expected & Status & Result & Time (s) & ATP calls \\
\midrule
\endfirsthead
\toprule
Proof & Expected & Status & Result & Time (s) & ATP calls \\
\midrule
\endhead
\bottomrule
\endlastfoot
\texttt{TST001{+}1} & correct & VerifiedGood & ok & 1.19 & 0 \\
\texttt{TST002{+}1} & correct & VerifiedGood & ok & 1.22 & 0 \\
\texttt{TST003{+}1} & correct & VerifiedGood & ok & 1.28 & 0 \\
\texttt{TST004{+}1} & correct & VerifiedGood & ok & 1.24 & 0 \\
\texttt{TST005{+}1} & correct & VerifiedGood & ok & 1.30 & 0 \\
\texttt{TST006{+}1} & correct & VerifiedGood & ok & 1.27 & 0 \\
\texttt{TST007{+}1} & correct & VerifiedGood & ok & 1.37 & 0 \\
\texttt{TST008{+}1} & incorrect & VerifiedBad & ok & 1.61 & 0 \\
\texttt{TST009{+}1} & incorrect & VerifiedBad & ok & 1.57 & 0 \\
\texttt{TST010{+}1} & incorrect & VerifiedBad & ok & 1.24 & 0 \\
\texttt{TST011{+}1} & incorrect & VerifiedBad & ok & 1.29 & 0 \\
\texttt{TST012{+}1} & incorrect & VerifiedBad & ok & 1.21 & 0 \\
\texttt{TST013{+}1} & incorrect & VerifiedBad & ok & 1.18 & 0 \\
\texttt{TST014{+}1} & incorrect & VerifiedBad & ok & 1.22 & 0 \\
\texttt{TST015{+}1} & incorrect & VerifiedBad & ok & 1.30 & 0 \\
\texttt{TST016{+}1} & incorrect & VerifiedBad & ok & 1.41 & 0 \\
\texttt{TST017{+}1} & incorrect & VerifiedBad & ok & 1.19 & 0 \\
\texttt{TST018{+}1} & incorrect & VerifiedBad & ok & 1.28 & 0 \\
\texttt{TST019{+}1} & incorrect & VerifiedBad & ok & 1.15 & 0 \\
\texttt{TST020{+}1} & incorrect & Unknown & abstain & 1.18 & 0 \\
\texttt{TST021{+}1} & incorrect & VerifiedBad & ok & 1.24 & 0 \\
\texttt{TST022{+}1} & incorrect & VerifiedBad & ok & 1.28 & 0 \\
\texttt{TST023{+}1} & incorrect & VerifiedBad & ok & 1.75 & 1 \\
\texttt{TST024{+}1} & incorrect & VerifiedBad & ok & 1.65 & 1 \\
\texttt{TST025{+}1} & incorrect & VerifiedBad & ok & 1.33 & 0 \\
\texttt{TST026{+}1} & incorrect & VerifiedBad & ok & 1.46 & 1 \\
\texttt{TST027{+}1} & incorrect & VerifiedBad & ok & 1.24 & 0 \\
\texttt{TST028{+}1} & incorrect & VerifiedBad & ok & 1.28 & 0 \\
\texttt{TST029{+}1} & incorrect & VerifiedBad & ok & 1.26 & 0 \\
\texttt{TST030{+}1} & incorrect & VerifiedBad & ok & 1.36 & 0 \\
\texttt{TST031{+}1} & incorrect & VerifiedBad & ok & 1.40 & 0 \\
\texttt{TST032{+}1} & incorrect & VerifiedBad & ok & 1.30 & 0 \\
\texttt{TST033{+}1} & incorrect & VerifiedBad & ok & 1.32 & 0 \\
\texttt{TST034{+}1} & incorrect & VerifiedBad & ok & 1.28 & 0 \\
\texttt{TST035{+}1} & incorrect & VerifiedBad & ok & 1.27 & 0 \\
\texttt{TST036{+}1} & incorrect & VerifiedBad & ok & 1.33 & 0 \\
\texttt{TST037{+}1} & incorrect & VerifiedBad & ok & 1.49 & 1 \\
\texttt{TST038{+}1} & incorrect & VerifiedBad & ok & 1.56 & 0 \\
\texttt{TST039{+}1} & incorrect & VerifiedBad & ok & 1.25 & 0 \\
\texttt{TST040{+}1} & incorrect & Error & ok & 1.24 & 0 \\
\texttt{TST041{+}1} & incorrect & VerifiedBad & ok & 1.29 & 0 \\
\texttt{TST042{+}1} & incorrect & VerifiedBad & ok & 1.21 & 0 \\
\texttt{TST043{+}1} & incorrect & VerifiedGood & wrong accept & 1.23 & 1 \\
\texttt{TST044{+}1} & incorrect & VerifiedBad & ok & 1.30 & 0 \\
\texttt{TST045{+}1} & incorrect & VerifiedBad & ok & 1.33 & 0 \\
\texttt{TST046{+}1} & incorrect & VerifiedBad & ok & 1.26 & 0 \\
\texttt{TST047{+}1} & correct & VerifiedGood & ok & 1.23 & 0 \\
\texttt{TST048{+}1} & correct & VerifiedGood & ok & 1.29 & 0 \\
\texttt{TST049{+}1} & correct & VerifiedGood & ok & 1.26 & 0 \\
\texttt{TST050{+}1} & correct & VerifiedGood & ok & 1.19 & 0 \\
\texttt{TST051{+}1} & correct & VerifiedGood & ok & 1.28 & 0 \\
\texttt{TST052{+}1} & correct & VerifiedGood & ok & 1.57 & 1 \\
\texttt{TST053{+}1} & incorrect & VerifiedBad & ok & 1.55 & 0 \\
\texttt{TST054{+}1} & incorrect & VerifiedBad & ok & 1.34 & 0 \\
\texttt{TST055{+}1} & incorrect & VerifiedBad & ok & 1.28 & 0 \\
\texttt{TST056{+}1} & incorrect & VerifiedBad & ok & 1.21 & 0 \\
\texttt{TST057{+}1} & incorrect & VerifiedBad & ok & 1.22 & 0 \\
\texttt{TST058{+}1} & incorrect & VerifiedBad & ok & 1.21 & 0 \\
\texttt{TST059{+}1} & incorrect & VerifiedBad & ok & 1.38 & 0 \\
\texttt{TST060{+}1} & incorrect & VerifiedBad & ok & 1.23 & 0 \\
\texttt{TST061{+}1} & incorrect & VerifiedBad & ok & 1.27 & 0 \\
\texttt{TST062{+}1} & incorrect & VerifiedBad & ok & 1.21 & 0 \\
\texttt{TST063{+}1} & incorrect & VerifiedBad & ok & 1.25 & 0 \\
\texttt{TST064{+}1} & incorrect & VerifiedBad & ok & 1.56 & 0 \\
\texttt{TST065{+}1} & incorrect & VerifiedBad & ok & 1.29 & 1 \\
\texttt{TST066{+}1} & incorrect & VerifiedBad & ok & 1.07 & 0 \\
\texttt{TST067{+}1} & incorrect & VerifiedBad & ok & 1.21 & 0 \\
\texttt{TST068{+}1} & edge & VerifiedGood & edge case & 1.39 & 0 \\
\texttt{TST069{+}1} & incorrect & VerifiedBad & ok & 1.22 & 0 \\
\texttt{TST070{+}1} & incorrect & VerifiedBad & ok & 1.28 & 0 \\
\texttt{TST071{+}1} & edge & VerifiedGood & edge case & 1.26 & 0 \\
\texttt{TST072{+}1} & incorrect & VerifiedBad & ok & 1.25 & 0 \\
\texttt{TST073{+}1} & edge & VerifiedGood & edge case & 1.23 & 0 \\
\texttt{TST074{+}1} & edge & VerifiedBad & edge case & 1.31 & 0 \\
\texttt{TST075{+}1} & incorrect & VerifiedBad & ok & 1.36 & 1 \\
\texttt{TST076{+}1} & edge & VerifiedBad & edge case & 1.32 & 1 \\
\end{longtable}

}

{\small
\begin{longtable}{@{}l l l l r r@{}}
\caption{Results for each proof of the Extra set (generated rule-level test proofs).}\label{tab:app-ext}\\
\toprule
Proof & Expected & Status & Result & Time (s) & ATP calls \\
\midrule
\endfirsthead
\toprule
Proof & Expected & Status & Result & Time (s) & ATP calls \\
\midrule
\endhead
\bottomrule
\endlastfoot
\texttt{EXT001{+}1} & correct & VerifiedGood & ok & 1.28 & 0 \\
\texttt{EXT002{+}1} & incorrect & VerifiedBad & ok & 1.31 & 0 \\
\texttt{EXT003{+}1} & incorrect & VerifiedBad & ok & 1.20 & 0 \\
\texttt{EXT004{+}1} & incorrect & VerifiedBad & ok & 1.20 & 0 \\
\texttt{EXT005{+}1} & incorrect & VerifiedBad & ok & 1.33 & 0 \\
\texttt{EXT006{+}1} & incorrect & VerifiedBad & ok & 1.20 & 0 \\
\texttt{EXT007{+}1} & correct & VerifiedGood & ok & 1.24 & 0 \\
\texttt{EXT008{+}1} & incorrect & VerifiedBad & ok & 1.50 & 0 \\
\texttt{EXT009{+}1} & edge & VerifiedBad & edge case & 1.47 & 0 \\
\texttt{EXT010{+}1} & correct & VerifiedGood & ok & 1.27 & 0 \\
\texttt{EXT011{+}1} & incorrect & VerifiedBad & ok & 1.30 & 0 \\
\texttt{EXT012{+}1} & incorrect & VerifiedBad & ok & 1.30 & 0 \\
\texttt{EXT013{+}1} & incorrect & VerifiedBad & ok & 1.38 & 0 \\
\texttt{EXT014{+}1} & incorrect & VerifiedBad & ok & 1.42 & 0 \\
\texttt{EXT015{+}1} & incorrect & VerifiedBad & ok & 1.35 & 0 \\
\texttt{EXT016{+}1} & incorrect & VerifiedBad & ok & 1.38 & 0 \\
\texttt{EXT017{+}1} & correct & VerifiedGood & ok & 1.21 & 0 \\
\texttt{EXT018{+}1} & incorrect & VerifiedBad & ok & 1.24 & 0 \\
\texttt{EXT019{+}1} & correct & VerifiedGood & ok & 1.32 & 0 \\
\texttt{EXT020{+}1} & incorrect & VerifiedBad & ok & 1.22 & 0 \\
\texttt{EXT021{+}1} & correct & VerifiedGood & ok & 1.30 & 0 \\
\texttt{EXT022{+}1} & incorrect & VerifiedBad & ok & 1.30 & 0 \\
\texttt{EXT023{+}1} & incorrect & VerifiedBad & ok & 1.57 & 0 \\
\texttt{EXT024{+}1} & incorrect & VerifiedBad & ok & 1.59 & 0 \\
\texttt{EXT025{+}1} & incorrect & VerifiedBad & ok & 1.27 & 0 \\
\texttt{EXT026{+}1} & correct & VerifiedGood & ok & 1.23 & 0 \\
\texttt{EXT027{+}1} & incorrect & VerifiedBad & ok & 1.42 & 0 \\
\texttt{EXT028{+}1} & incorrect & VerifiedBad & ok & 1.24 & 0 \\
\texttt{EXT029{+}1} & incorrect & VerifiedBad & ok & 1.22 & 0 \\
\texttt{EXT030{+}1} & incorrect & VerifiedBad & ok & 1.25 & 0 \\
\texttt{EXT031{+}1} & correct & VerifiedGood & ok & 1.30 & 0 \\
\texttt{EXT032{+}1} & incorrect & VerifiedBad & ok & 1.20 & 0 \\
\texttt{EXT033{+}1} & incorrect & VerifiedBad & ok & 1.27 & 0 \\
\texttt{EXT034{+}1} & correct & VerifiedGood & ok & 1.32 & 0 \\
\texttt{EXT035{+}1} & incorrect & VerifiedBad & ok & 1.50 & 1 \\
\texttt{EXT036{+}1} & incorrect & VerifiedBad & ok & 1.34 & 1 \\
\texttt{EXT037{+}1} & incorrect & VerifiedBad & ok & 1.43 & 1 \\
\texttt{EXT038{+}1} & edge & VerifiedGood & edge case & 1.56 & 1 \\
\texttt{EXT039{+}1} & correct & VerifiedGood & ok & 1.24 & 0 \\
\texttt{EXT040{+}1} & incorrect & VerifiedBad & ok & 1.25 & 1 \\
\texttt{EXT041{+}1} & incorrect & VerifiedBad & ok & 1.29 & 1 \\
\texttt{EXT042{+}1} & correct & VerifiedGood & ok & 1.26 & 0 \\
\texttt{EXT043{+}1} & incorrect & VerifiedBad & ok & 1.32 & 1 \\
\texttt{EXT044{+}1} & correct & VerifiedGood & ok & 1.26 & 0 \\
\texttt{EXT045{+}1} & incorrect & VerifiedBad & ok & 1.36 & 1 \\
\texttt{EXT046{+}1} & incorrect & VerifiedBad & ok & 1.41 & 1 \\
\texttt{EXT047{+}1} & incorrect & VerifiedBad & ok & 1.30 & 1 \\
\texttt{EXT048{+}1} & incorrect & VerifiedBad & ok & 1.20 & 0 \\
\texttt{EXT049{+}1} & correct & VerifiedGood & ok & 1.38 & 0 \\
\texttt{EXT050{+}1} & incorrect & VerifiedBad & ok & 1.70 & 0 \\
\texttt{EXT051{+}1} & incorrect & VerifiedBad & ok & 1.32 & 0 \\
\texttt{EXT052{+}1} & incorrect & VerifiedBad & ok & 1.58 & 0 \\
\texttt{EXT053{+}1} & correct & VerifiedGood & ok & 1.64 & 0 \\
\texttt{EXT054{+}1} & incorrect & VerifiedBad & ok & 1.66 & 0 \\
\texttt{EXT055{+}1} & correct & VerifiedGood & ok & 1.38 & 0 \\
\texttt{EXT056{+}1} & incorrect & VerifiedBad & ok & 1.50 & 0 \\
\texttt{EXT057{+}1} & correct & VerifiedGood & ok & 1.33 & 0 \\
\texttt{EXT058{+}1} & correct & VerifiedGood & ok & 1.39 & 0 \\
\texttt{EXT059{+}1} & incorrect & VerifiedBad & ok & 1.56 & 0 \\
\texttt{EXT060{+}1} & incorrect & VerifiedBad & ok & 1.66 & 0 \\
\texttt{EXT061{+}1} & correct & VerifiedGood & ok & 1.38 & 0 \\
\texttt{EXT062{+}1} & incorrect & VerifiedBad & ok & 1.32 & 0 \\
\texttt{EXT063{+}1} & incorrect & VerifiedBad & ok & 1.35 & 0 \\
\texttt{EXT064{+}1} & incorrect & VerifiedBad & ok & 1.25 & 0 \\
\texttt{EXT065{+}1} & correct & VerifiedGood & ok & 1.24 & 0 \\
\texttt{EXT066{+}1} & incorrect & VerifiedBad & ok & 1.65 & 0 \\
\texttt{EXT067{+}1} & correct & VerifiedGood & ok & 1.56 & 0 \\
\texttt{EXT068{+}1} & incorrect & VerifiedBad & ok & 1.59 & 0 \\
\texttt{EXT069{+}1} & incorrect & VerifiedBad & ok & 1.61 & 0 \\
\texttt{EXT070{+}1} & edge & VerifiedBad & edge case & 1.30 & 0 \\
\texttt{EXT071{+}1} & correct & VerifiedGood & ok & 1.27 & 0 \\
\texttt{EXT072{+}1} & incorrect & VerifiedBad & ok & 1.31 & 0 \\
\texttt{EXT073{+}1} & incorrect & VerifiedBad & ok & 1.24 & 0 \\
\texttt{EXT074{+}1} & correct & VerifiedGood & ok & 1.25 & 0 \\
\texttt{EXT075{+}1} & incorrect & VerifiedBad & ok & 1.29 & 0 \\
\texttt{EXT076{+}1} & incorrect & VerifiedBad & ok & 1.30 & 0 \\
\texttt{EXT077{+}1} & correct & VerifiedGood & ok & 1.31 & 0 \\
\texttt{EXT078{+}1} & incorrect & VerifiedBad & ok & 1.34 & 0 \\
\texttt{EXT079{+}1} & incorrect & VerifiedBad & ok & 1.22 & 0 \\
\texttt{EXT080{+}1} & correct & VerifiedGood & ok & 1.25 & 0 \\
\texttt{EXT081{+}1} & incorrect & VerifiedBad & ok & 1.43 & 0 \\
\texttt{EXT082{+}1} & incorrect & VerifiedBad & ok & 1.43 & 0 \\
\texttt{EXT083{+}1} & correct & VerifiedGood & ok & 1.20 & 0 \\
\texttt{EXT084{+}1} & incorrect & VerifiedBad & ok & 1.28 & 0 \\
\texttt{EXT085{+}1} & incorrect & VerifiedBad & ok & 1.18 & 0 \\
\texttt{EXT086{+}1} & correct & VerifiedGood & ok & 1.19 & 0 \\
\texttt{EXT087{+}1} & incorrect & VerifiedBad & ok & 1.23 & 0 \\
\texttt{EXT088{+}1} & incorrect & VerifiedBad & ok & 1.42 & 0 \\
\texttt{EXT089{+}1} & correct & VerifiedGood & ok & 1.34 & 0 \\
\texttt{EXT090{+}1} & incorrect & VerifiedBad & ok & 1.48 & 1 \\
\texttt{EXT091{+}1} & correct & VerifiedGood & ok & 1.45 & 0 \\
\texttt{EXT092{+}1} & incorrect & VerifiedBad & ok & 1.36 & 0 \\
\texttt{EXT093{+}1} & correct & VerifiedGood & ok & 1.34 & 0 \\
\texttt{EXT094{+}1} & incorrect & VerifiedBad & ok & 1.25 & 0 \\
\texttt{EXT095{+}1} & incorrect & VerifiedBad & ok & 1.53 & 0 \\
\texttt{EXT096{+}1} & correct & VerifiedGood & ok & 2.64 & 0 \\
\texttt{EXT097{+}1} & edge & VerifiedBad & edge case & 1.45 & 0 \\
\texttt{EXT098{+}1} & edge & VerifiedBad & edge case & 1.45 & 0 \\
\texttt{EXT099{+}1} & correct & VerifiedGood & ok & 1.35 & 0 \\
\texttt{EXT100{+}1} & incorrect & VerifiedBad & ok & 1.35 & 0 \\
\texttt{EXT101{+}1} & correct & VerifiedGood & ok & 1.35 & 0 \\
\texttt{EXT102{+}1} & incorrect & VerifiedBad & ok & 1.52 & 0 \\
\texttt{EXT103{+}1} & edge & VerifiedBad & edge case & 1.28 & 0 \\
\texttt{EXT104{+}1} & edge & VerifiedGood & edge case & 1.39 & 1 \\
\texttt{EXT105{+}1} & edge & VerifiedBad & edge case & 1.25 & 0 \\
\texttt{EXT106{+}1} & correct & VerifiedGood & ok & 1.66 & 0 \\
\texttt{EXT107{+}1} & correct & VerifiedGood & ok & 1.37 & 0 \\
\texttt{EXT108{+}1} & correct & VerifiedGood & ok & 1.63 & 0 \\
\texttt{EXT109{+}1} & incorrect & VerifiedBad & ok & 1.78 & 1 \\
\end{longtable}

}
